\chapter{Lernziele Kompression}
\begin{todolist}
\item Ich kann Beispiele aus der Geschichte angeben, die aufzeigen, dass Kompression schon früh ein Bedürfnis war.
\item Ich kann den Unterschied zwischen verlustbehafteten und verlustfreien Kompressionen erklären und Beispiele dafür angeben.
\item Ich kann an einem Beispiel erklären, wie Verweise auf bereits vorgekommene Wörter einen Text verkürzen (Wörterbuchkompression).
\item Ich kann einen kurzen Text mit Verweisen (Abstand, Länge) kodieren und eine solche Folge dekodieren (LZ77).
\item Ich weiss, was eine präfixfreie Kodierung ist, und wieso diese Eigenschaft wichtig ist. Ich kann prüfen, ob eine Kodierung präfixfrei ist.
\item Ich kann mit einem Baumdiagramm eine präfixfreie Kodierung finden.
\item Ich kann die absoluten und relativen Häufigkeiten von Zeichen in einem (kurzen) Text bestimmen.
\item Ich kann berechnen, wie viele Bits ein Text bei einer Kodierung mit fester Codelänge braucht.
\item Ich kann eine präfixfreie Kodierung mit der Methode der maximalen Balance für einen Text finden, von dem die relativen Häufigkeiten der Zeichen gegeben sind.
\item Ich kann ein Beispiel für eine relative Häufigkeit von Buchstaben angeben, bei welchem die Kodierung der maximalen Balance einen selteneren Buchstaben kürzer kodiert als einen häufigeren.
\item Ich kann eine Huffman-Kodierung für einen Text finden, von dem die relativen Häufigkeiten der Zeichen gegeben sind.
\item Ich kann ausgehend von der Kodierung und den absoluten Häufigkeiten der vorkommenden Zeichen berechnen, wie viele Bits für die Speicherung eines Textes nötig sind.
\item Ich kann die mittlere Codelänge (Bit pro Zeichen) einer Kodierung berechnen und damit die Methode der maximalen Balance mit der Huffman-Kodierung vergleichen.
\item Ich kann erklären, wie ZIP (DEFLATE) LZ77 und die Huffman-Kodierung kombiniert, und die Anzahl Bits eines kurzen Texts für alle Verfahren vergleichen.
\item Ich kann bei gegebener relativer Häufigkeit der Zeichen einen (kurzen) Text arithmetisch kodieren.
\item Ich kann bei gegebener relativer Häufigkeit der Zeichen eine arithmetische Kodierung dekomprimieren, die als Intervall oder als Zahl mit Textlänge gegeben ist.
\end{todolist}